% Standalone problem document, generated from the adjacent README.md.
% Compile with XeLaTeX. Mathematical content is maintained in Markdown.
\documentclass[11pt,a4paper]{article}
\usepackage[fontsize=10.5pt]{scrextend}
\usepackage[margin=22mm,headheight=15pt,headsep=9mm,footskip=11mm]{geometry}
\usepackage{fontspec}
\setmainfont{texgyrepagella}[Extension=.otf,UprightFont=*-regular,BoldFont=*-bold,ItalicFont=*-italic,BoldItalicFont=*-bolditalic]
\setsansfont{texgyreheros}[Extension=.otf,UprightFont=*-regular,BoldFont=*-bold,ItalicFont=*-italic,BoldItalicFont=*-bolditalic]
\setmonofont{texgyrecursor}[Extension=.otf,UprightFont=*-regular,BoldFont=*-bold,ItalicFont=*-italic,BoldItalicFont=*-bolditalic,Scale=MatchLowercase]
\usepackage{amsmath,amssymb}
\usepackage{unicode-math}
\setmathfont{texgyrepagella-math.otf}
\usepackage{xcolor,microtype,enumitem,needspace,fancyhdr}
\usepackage{bookmark}
\definecolor{ink}{HTML}{17344B}
\definecolor{accent}{HTML}{197D80}
\definecolor{muted}{HTML}{596773}
\hypersetup{colorlinks=true,linkcolor=accent,urlcolor=accent,pdftitle={IE-12 - Near-quadratic
solution cost at a prescribed backward
error},pdfauthor={Open Problems in Numerical Linear Algebra}}
\urlstyle{same}
\setlength{\parindent}{0pt}
\setlength{\parskip}{5pt plus 1pt minus 1pt}
\linespread{1.04}
\setlength{\emergencystretch}{3em}
\widowpenalty=10000
\clubpenalty=10000
\displaywidowpenalty=10000
\allowdisplaybreaks[1]
\setlist{nosep,leftmargin=1.5em}
\providecommand{\tightlist}{\setlength{\itemsep}{0pt}\setlength{\parskip}{0pt}}
\providecommand{\pandocbounded}[1]{#1}
\pagestyle{fancy}
\fancyhf{}
\fancyhead[L]{\sffamily\footnotesize\color{muted}OPEN PROBLEMS IN NUMERICAL LINEAR ALGEBRA}
\fancyhead[R]{\sffamily\bfseries\footnotesize\color{accent}IE-12}
\fancyfoot[L]{\parbox[b]{0.9\textwidth}{\sffamily\fontsize{8}{10}\selectfont\color{muted}Editorial ratings; a bounded search cannot certify that a problem remains unsolved.\\Literature check: 2026-09-12}}
\fancyfoot[R]{\sffamily\footnotesize\color{muted}\thepage}
\renewcommand{\headrulewidth}{0.3pt}
\renewcommand{\headrule}{\hbox to\headwidth{\color{accent}\leaders\hrule height \headrulewidth\hfill}}
\makeatletter
\renewcommand{\section}{\Needspace{5\baselineskip}\@startsection{section}{1}{0pt}{12pt plus 2pt minus 2pt}{4pt}{\sffamily\bfseries\large\color{ink}}}
\renewcommand{\subsection}{\Needspace{4\baselineskip}\@startsection{subsection}{2}{0pt}{9pt plus 1pt minus 1pt}{3pt}{\sffamily\bfseries\normalsize\color{ink}}}
\makeatother
\setcounter{secnumdepth}{0}
\begin{document}
{\sffamily\small\color{muted}Linear systems and elimination\par}
\vspace{3pt}
{\raggedright\hyphenpenalty=10000\exhyphenpenalty=10000\sffamily\bfseries\fontsize{21}{24}\selectfont\color{ink}Near-quadratic
solution cost at a prescribed backward error\par}
\vspace{7pt}
{\sffamily\small\textbf{Difficulty:} challenging\quad\textbf{Importance:} broadly
interesting\par}
{\sffamily\small\bfseries\color{accent}Status: Solved\par}
\vspace{4pt}
{\color{accent}\hrule height 0.8pt}
\vspace{6pt}
\textbf{Rating rationale (historical):} Challenging reflects the
remaining removal of a dimension-dependent logarithm from an established
general algorithm; broad impact follows from a condition-independent
cost bound for arbitrary linear systems.

\subsection{Resolution: affirmative, 12 September
2026}\label{resolution-affirmative-12-september-2026}

\textbf{Author:} Sidney Holden, Center for Computational Biology,
Flatiron Institute, Simons Foundation.
\href{https://github.com/ajt60gaibb/OpenProblemsInNLA/blob/main/references/holden-ie12-2026-09-12/README.md}{Verified
affiliation and submission record}.

\href{https://github.com/ajt60gaibb/OpenProblemsInNLA/blob/main/linear-systems-and-elimination/IE-12/solution.pdf}{Theorem
1, proved in Sections 2--6} gives a randomized exact-real algorithm with
deterministic worst-case cost \(O(n^2\varepsilon^{-3})\) and \(O(n^2)\)
scalar storage. For every original input it always returns a nonzero
vector and achieves the original A-only backward error at most
\(5\varepsilon/8\) with probability greater than \(0.997\). Thus it
settles the complete target with \(q=3\), uniformly in dimension,
conditioning and right-hand side.
\href{https://github.com/ajt60gaibb/OpenProblemsInNLA/blob/main/linear-systems-and-elimination/IE-12/solution.tex}{Proof
source}.

The argument passed a separate
\href{https://github.com/ajt60gaibb/OpenProblemsInNLA/blob/main/references/holden-ie12-2026-09-12/independent-review.md}{independent
Codex AI-agent audit}, including exact component-test reproduction. This
is informal automated review, not external human peer review or formal
verification. The supplied draft was prepared with ChatGPT assistance;
no Lean verification was performed.

Entrywise randomized rounding and weighted-pattern matrix products
cancel the logarithmic iteration overhead in total arithmetic cost. This
does not assert dimension-independent black-box convergence,
finite-precision stability, bit complexity or practical speed. The
original target and prior-source attribution below are retained; the
earlier status discussion is historical.

\subsection{Problem statement}\label{problem-statement}

Does there exist a randomized algorithm and absolute constants \(C,q>0\)
such that, for every \(n\ge1\), nonsingular
\(A\in\mathbb R^{n\times n}\) with \(\|A\|_2=1\), \(b\ne0\), and
\(0<\varepsilon<1/2\), it returns \(x\ne0\) satisfying

\[\Pr\left\{\frac{\|Ax-b\|_2}{\|A\|_2\|x\|_2}
\le\varepsilon\right\}\ge0.99\]

using at most \(C n^2\varepsilon^{-q}\) operations? Use an exact-real
arithmetic model with scalar arithmetic, square roots, comparisons, and
independent standard Gaussian or random-bit draws at unit cost. All
input entries may be accessed; the norm normalization is an input
promise. The cost constants must be independent of the condition number
and of \(n,b,\varepsilon\). The fraction is normwise backward error with
perturbations allowed in \(A\) only. A deterministic algorithm meeting
the bound would also answer positively.

\subsection{References and status}\label{references-and-status}

M. Dereziński, Y. Nakatsukasa, and E. Rebrova,
\href{https://arxiv.org/html/2604.16075v2}{\emph{Towards Universal
Convergence of Backward Error in Linear System Solvers}} (22 May 2026),
§§1, 3.1, 5.2, 7, poses the cost question and specifies the
backward-error metric. Corollary 17 gives \(O(n^2/\sqrt\varepsilon)\)
for PSD systems. The revised paper's Corollary 25 proves the
general-system bound \(O(n^2\log(n/\delta)/\varepsilon)\) with failure
probability \(\delta\); its §7 asks for convergence independent of
dimension. The remaining target here removes the \(\log n\) overhead,
which cannot be hidden in a constant depending only on \(\varepsilon\).
The model and success probability above make this explicit. Searches on
2026-09-08 for ``universal backward error 2026 linear systems'',
``MINBERR smoothed analysis'', and ``quadratic backward error linear
solver'' found no resolution of the displayed bound. The earlier
version's lack of a proved general-system rate is no longer current.

\subsection{Audit update --- 2026-09-10}\label{audit-update-2026-09-10}

Rechecked \href{https://arxiv.org/html/2604.16075v2}{version 2},
Corollaries 17 and 25 and §7. Positive definite inputs satisfy the
displayed cost by Corollary 17; the general bound still contains
\(\log n\). This warrants partial status and a difficulty change from
extreme to challenging for the narrowed remaining gap. Searches for
universal backward-error convergence and MINBERR follow-ups found no
dimension-independent general bound.
\end{document}
